\section{Results II: the three tiers}
\label{sec:results-models}

All results in this section are fitted on \code{dataset\_v1.1} with unloading measured at the day of the scan.

\subsection{Tier 1: the duration--response}
\label{sec:results-tier1}

\paragraph{The shape of the curve} The three duration forms fit the 346 rows of subset A almost equally well
(\cref{tab:forms}): the logarithmic, saturating and spline forms lie within 1.3 AIC points of one another, far
inside the 4-point margin the headline rule required for another form to displace the saturating one. The data
therefore establish that loss accumulates along a curved path, but not \emph{which} curve. The spline shows no shape
the parametric forms miss ($\Delta\mathrm{AIC} = -0.2$). By the declared rule the saturating form is the headline.

\begin{table}[htbp]
\centering
\caption[Tier 1: the three duration forms]{Tier 1, subset A (346 rows, 32 campaigns): the three duration forms.
$\Delta$AIC is relative to the headline saturating form; $k$ counts fixed effects, the four variance components
and, for the saturating form, $\tau$. Intervals are 95\% cluster-robust ($t_{31}$).}
\label{tab:forms}
\small
\begin{tabularx}{\linewidth}{@{}L{3.6cm} r r r Y@{}}
\toprule
Form & AIC & $\Delta$AIC & $k$ & Duration term \\
\midrule
\inputtable{generated/tab_tier1_forms.tex}
\bottomrule
\end{tabularx}
\end{table}

\paragraph{How fast and how far} Under the logarithmic form, muscle size falls by 2.99\pp{} for every doubling of
days (95\% CI 2.53 to 3.44). Under the saturating form, the time constant is 60 days---the profile likelihood turns
over inside the declared grid---and the eventual loss is $-17.3\%$ ($-20.9$ to $-13.7$) for the reference scenario.
By day 119, the longest observation, the curve has covered 86\% of the way to that plateau, so the plateau is an
estimate at the edge of the data rather than far beyond it. \Cref{fig:duration} shows the curve with the
control-arm observations, and \cref{tab:curve} its values at selected days.

\begin{figure}[htbp]
\centering
\includegraphics[width=\linewidth]{fig_duration}
\caption[Duration--response: control-arm observations and the tier-1 curve]{Duration--response. Points are the
246 control-arm rows of subset A (30 campaigns), placed at the day of the scan and coloured by functional class. The
line is the tier-1 saturating curve for the reference scenario (a control arm, knee extensors, CT cross-sectional
area, not a composite) with its 95\% cluster-robust band. The scatter around the curve is dominated by which muscle
was measured: antigravity extensors lie mostly below it, flexors and mixed muscles above.}
\label{fig:duration}
\end{figure}

\begin{table}[htbp]
\centering
\caption[Tier 1: the fitted curve at selected days]{The tier-1 saturating curve at selected days for the reference
scenario, with 95\% cluster-robust confidence intervals.}
\label{tab:curve}
\footnotesize
\setlength{\tabcolsep}{4pt}
\begin{tabular}{@{}l cccccc@{}}
\toprule
\inputtable{generated/tab_curve.tex}
\bottomrule
\end{tabular}
\end{table}

\paragraph{Where the variance sits} \Cref{tab:variance} decomposes the variance. In subset A, 54\% is residual,
30\% lies \emph{between muscles within a campaign}, 11\% between campaigns and 4\% between papers reporting the same
campaign. The 30\% is the quantitative case for the second claim of the project: once duration is accounted for,
which muscle is measured matters more than which campaign measured it.

\begin{table}[htbp]
\centering
\caption[Tier 1: variance components]{Tier 1 variance components (squared percentage points) and their share of
the total.}
\label{tab:variance}
\small
\begin{tabular}{@{}lcccc@{}}
\toprule
Model & Campaign & Study within campaign & Muscle within campaign & Residual \\
\midrule
\inputtable{generated/tab_variance.tex}
\bottomrule
\end{tabular}
\end{table}

\paragraph{Other terms} In subset A, countermeasure arms lose 3.43\pp{} less than control arms (95\% CI 1.68 to
5.18, $p < 0.001$). Relative to CT cross-sectional area, DXA lean mass shows 4.75\pp{} less loss (0.05 to 9.45,
$p = 0.048$) and ultrasound thickness 5.95\pp{} less (0.61 to 11.29, $p = 0.030$), while MRI volume differs by only
0.49\pp{} ($p = 0.72$); a curve drawn for MRI volume, the dominant measurement, would therefore sit about half a
percentage point above the reference curve. The composite flag and the whole-segment family carry no detectable
effect. All coefficients are listed in \cref{app:fixed}.

\subsection{Tier 1: muscle-specific vulnerability}
\label{sec:results-ranking}

The ranking model, fitted on the 304 named-muscle rows of subset B (25 campaigns), orders the muscle families by
their fitted change at day 60 for a control arm (\cref{tab:ranking,fig:ranking}). The plantar flexors lose most,
$-15.8\%$ ($-18.9$ to $-12.7$), significantly more than the knee extensors ($-5.69$\pp, $p = 0.001$). The knee
extensors, knee flexors and dorsiflexors form a middle group near $-10\%$; the hip adductors, extensors and rotators
lose significantly less than the knee extensors, the hip rotators least of all ($-1.9\%$).

\begin{table}[htbp]
\centering
\caption[Tier 1: muscle-family ranking at day 60]{Tier 1 muscle-family ranking, subset B (304 rows, 25 campaigns):
fitted change at day 60 for a control arm, and contrast with the reference family (knee extensors). Intervals are
95\% cluster-robust ($t_{24}$). Hip abductors and hip rotators each rest on a single campaign.}
\label{tab:ranking}
\small
\begin{tabular*}{\linewidth}{@{\extracolsep{\fill}}l r r c c r@{}}
\toprule
Family & Rows & Campaigns & \makecell{Change at day 60,\\ \% (95\% CI)} & \makecell{vs.\ knee extensors,\\ pp (95\% CI)} & $p$ \\
\midrule
\inputtable{generated/tab_ranking.tex}
\bottomrule
\end{tabular*}
\end{table}

\begin{figure}[htbp]
\centering
\includegraphics[width=\linewidth]{fig_ranking}
\caption[Muscle-family ranking at day 60]{Muscle-family ranking at day 60 for a control arm, with 95\%
cluster-robust intervals, from the tier-1 ranking model on subset B. The dashed line marks the reference family.
Rows and campaigns are shown for each family because two families rest on a single campaign.}
\label{fig:ranking}
\end{figure}

\paragraph{The antigravity contrast} Because every contrast in \cref{tab:ranking} inherits the uncertainty of the
reference, the comparison that carries the physiological claim is reported on its own: plantar flexors against
dorsiflexors, the antagonist pair at the ankle, is $-5.60$\pp{} (95\% CI $-9.13$ to $-2.07$, $p = 0.003$), computed
from the coefficient difference so that no choice of reference can change it. The ordering survives adjustment for
duration, arm and modality.

\paragraph{Countermeasures} In the ranking model, countermeasure arms lose 3.73\pp{} less than control arms (95\% CI
1.54 to 5.93, $p = 0.002$). This is the first countermeasure effect the project has estimated rather than described;
it pools every countermeasure type and says nothing about which works best.

\subsection{Sensitivity analysis S6: does the weighting matter?}
\label{sec:results-s6}

The primary analysis weights rows by the number of participants analysed. The classical alternative, inverse-variance
weighting, needs a dispersion \emph{of the change}, which only 161 of the 346 rows in subset A, from 8 of the 32
campaigns, carry; most papers print the spread of the baseline or nothing. \Cref{tab:s6} compares the two schemes on
the same rows. The restriction to rows with a change dispersion moves the duration coefficient by 0.3\pp{} and the
weighting scheme by another 0.9\pp, with overlapping intervals throughout: on this literature the scheme barely
matters. The finding that belongs in the methods is the other one: inverse-variance weighting is unavailable as a
primary scheme, because only eight campaigns carry what it needs.

\begin{table}[htbp]
\centering
\caption[Sensitivity analysis S6]{Sensitivity analysis S6. The duration coefficient is the saturating term of the
tier-1 model; the error is out-of-campaign and campaign-weighted. From \file{results/sensitivity.md}.}
\label{tab:s6}
\small
\begin{tabularx}{\linewidth}{@{}Y l r r c r@{}}
\toprule
Analysis & Weights & Rows & Campaigns & Duration coefficient (95\% CI) & MAE (pp) \\
\midrule
Primary: every row & participants & 346 & 32 & $-14.25$ ($-16.40$ to $-12.11$) & 3.34 \\
Rows with a dispersion of the change & participants & 161 & 8 & $-14.51$ ($-15.61$ to $-13.42$) & 2.67 \\
S6: the same rows & inverse variance & 161 & 8 & $-13.60$ ($-16.69$ to $-10.50$) & 2.82 \\
\bottomrule
\end{tabularx}
\end{table}

\subsection{The duration-only baseline}
\label{sec:results-baseline}

Fitted inside the leave-one-campaign-out loop, the logarithmic curve is the best of the three baseline forms, with
an out-of-campaign MAE of 3.16\pp{} (95\% CI 2.44 to 3.93); the saturating form follows closely (3.19\pp) and the
straight line is worst (3.40\pp), which is the evidence that the relationship is curved rather than assumed to be
(\cref{tab:models}). All three forms already meet the 4\pp{} target before any model family is fitted. Their pooled
$R^2$ is low (0.07 to 0.14): duration alone predicts the average loss at a given day, but explains little of the
row-to-row variation, because that variation is mostly which muscle was measured. Averaged over all muscles in
subset A, the fitted saturating baseline has a time constant of 35 days and an eventual loss of $-13.0\%$; the
whole-segment and composite rows pull that average towards the less affected tissue, which is why the tier-1 model
with muscle terms, not the baseline, answers the question for a particular muscle.

\subsection{Tier 2: the comparative machine-learning families}
\label{sec:results-tier2}

\begin{keybox}{Tier 2 is a null result}
No family beats the duration-only curve (\cref{tab:models,fig:models}). The best, a random forest, scores 3.18\pp{}
against the curve's 3.16\pp{}; the others are between 3.7\% and 5.9\% worse. The plan set the bar at a 15\% relative
improvement and pre-committed to reporting the outcome either way. This is rung~B of the project's fallback ladder.
A pretrained tabular neural network, TabPFN, added afterwards as a fifth family, does not beat the curve either
(3.22\pp, 2.1\% worse).
\end{keybox}

\begin{table}[htbp]
\centering
\caption[Tier 2: out-of-campaign error of the baseline and the four families]{Out-of-campaign error of the
duration-only baseline in its three forms and of the four tier-2 families, subset A, 32 leave-one-campaign-out
folds. MAE and RMSE are campaign-weighted; the interval is a campaign bootstrap; ``vs.\ baseline'' is the relative
MAE change against the logarithmic curve (negative = worse); $R^2$ is pooled over out-of-fold predictions. The last
column names the campaign with the largest error and that error. TabPFN was added after the four families' result
was known and was run in its own environment (\cref{sec:tier2}).}
\label{tab:models}
\footnotesize
\begin{tabular*}{\linewidth}{@{\extracolsep{\fill}}l c c c c c l@{}}
\toprule
Model & \makecell{MAE\\ (pp)} & 95\% CI & RMSE & \makecell{vs.\\ baseline} & $R^2$ & Worst campaign (MAE) \\
\midrule
\inputtable{generated/tab_models.tex}
\bottomrule
\end{tabular*}
\end{table}

\begin{figure}[htbp]
\centering
\includegraphics[width=\linewidth]{fig_models}
\caption[Tier 2 against the duration curve]{Tier 2 against the duration curve. (a)~Out-of-campaign MAE with 95\%
campaign-bootstrap intervals; the dashed line is the logarithmic curve. (b)~Pooled out-of-fold $R^2$. Blue: the
four declared families; orange: TabPFN, added post hoc. The flexible families explain two to three times more
variance than the curve, yet do not convert it into lower error.}
\label{fig:models}
\end{figure}

The informative part is panel~b of \cref{fig:models}. Pooled $R^2$ rises from 0.14 for the curve to 0.27--0.36
once muscle identity enters: the flexible families do find real structure, but they do not convert it into a lower
absolute error, because the residual is dominated by between-campaign differences that no column in the dataset
explains. The worst folds show where: a 21-day campaign with site-resolved thigh measurements (HDBR-21, Liphardt;
7.8--7.9\pp) for three families, and WISE-2005, 60 days of bed rest in women only, for ridge regression (8.6\pp).

\paragraph{The pretrained network} TabPFN, the post-hoc fifth family, scores 3.22\pp{} (95\% CI 2.51 to 3.97), 2.1\%
worse than the curve and second of the five families behind the random forest. Its pooled $R^2$ of 0.39 is the highest
of any tier-2 model: a network pretrained on millions of synthetic tables, conditioning on the training campaigns in
context, finds the muscle structure as the tuned families do, and like them does not turn it into lower error in a
campaign it has not seen. Its worst campaigns are the same ones: HDBR-21 (Liphardt; 7.56\pp) and WISE-2005 (7.34\pp).
An earlier run of the same family, on \code{dataset\_v1.0} and the campaigns' planned lengths, had scored 3.25\pp{} with
a pooled $R^2$ of 0.03, behaving like a duration-only curve; on the day of the scan it behaves like the other families.
Either way the conclusion is the same, and it now rests on five families of very different kinds, from penalised
regression to a pretrained transformer.

\paragraph{Importance} SHAP values on the random forest identify two stable features (\cref{tab:importance}): the
duration term reaches the top three in 29 of 32 folds (91\%) and the plantar-flexor indicator in 28 (88\%). The third
place alternates between the countermeasure arm (59\%) and the composite flag (53\%), so the declared criterion---the
same top three in 80\% of folds---is missed. The flexible models lean on exactly the two variables the tier-1 model
estimates with intervals, which is the consistency the importance analysis was meant to check; it is not evidence of
cause.

\begin{table}[htbp]
\centering
\caption[Tier 2: stability of feature importance]{Stability of SHAP importance for the random forest across the 32
folds: how often each feature reaches the top three, and its mean absolute SHAP value. Features never in the top three
are omitted.}
\label{tab:importance}
\small
\begin{tabular}{@{}lccc@{}}
\toprule
Feature & Folds in top three & Share & Mean $|$SHAP$|$ \\
\midrule
\inputtable{generated/tab_importance.tex}
\bottomrule
\end{tabular}
\end{table}

\subsection{Tier 3: the probabilistic forecast}
\label{sec:results-tier3}

The first run sent 430 requests, the largest of 26,390 input tokens. \Cref{tab:forecast} reports both arms against
their references on identical metrics.

\begin{table}[htbp]
\centering
\caption[Tier 3: forecast accuracy and calibration]{Tier 3 against its references. MAE and CRPS in percentage
points, campaign-weighted, with a campaign-bootstrap interval for the MAE; the log score is per row (lower is better);
coverage and mean width (pp) of the most probable 50\% and 80\% ranges. $^\dagger$ the arm's reference.}
\label{tab:forecast}
\small
\begin{tabularx}{\linewidth}{@{}Y c c c c c c c@{}}
\toprule
Model & MAE & 95\% CI & CRPS & Log score & 50\% range & 80\% range & $R^2$ \\
\midrule
\inputtable{generated/tab_forecast.tex}
\bottomrule
\end{tabularx}
\end{table}

\paragraph{Without history} \Jev's error is 2.71\pp{} against the curve's 3.13\pp.\footnote{The curve scores
3.13\pp{} here and 3.16\pp{} in \cref{tab:models} because here its point is the mean of its forecast distribution, the
same definition used for the model's point.} Paired by campaign, the model's error is lower by 0.42\pp{} (95\% CI
0.13 to 0.73), and it is closer in 20 of the 32 campaigns (\cref{fig:jevcampaigns}). This is the first advantage over
the duration curve in the project whose paired interval excludes zero, and the pooled $R^2$ of 0.43 is about three
times the curve's (\cref{fig:jevscatter}). Against the declared bars, however, the MAE improvement is 13.3\%, short
of the 15\% required; the CRPS improvement, 17.7\%, meets its bar; and the 80\% ranges contain the truth only 62\% of
the time, so they may not be called 80\% ranges. The log score says the same from the other side: it is worse than
the curve's (2.84 against 2.29) despite the better CRPS, because when the model is wrong it has placed almost no
probability where the truth landed.

\paragraph{With history} Forecasting the change since the previous scan on 84 rows from five campaigns, \Jev{}
reaches 2.84\pp{} against 3.21\pp{} for the last scan plus the curve's step, a paired gain of 0.37\pp{} (0.03 to
0.75) in four of the five campaigns. The MAE improvement (11.5\%) misses its bar, the CRPS is 2.2\% worse than the
reference, and the ranges are badly overconfident: 2.8\pp{} wide, they contain the truth 38\% of the time. With five
campaigns, every interval on this arm rests on a bootstrap over five values and is indicative only.

\begin{figure}[htbp]
\centering
\includegraphics[width=\linewidth]{fig_jev_campaigns}
\caption[Tier 3: gain by held-out campaign]{Tier 3 without history: the duration curve's campaign MAE minus \Jev's,
for each held-out campaign (positive = \Jev{} closer). The number beside each bar is the campaign's rows. More than
half of the total gain comes from three long, multi-muscle MRI campaigns---NASA SPRINT, Berlin BBR2-2 and MEDES
LTBR---which are also among the most published campaigns in the field.}
\label{fig:jevcampaigns}
\end{figure}

\begin{figure}[htbp]
\centering
\includegraphics[width=\linewidth]{fig_jev_scatter}
\caption[Tier 3: forecast against observed value]{Forecast against observed change for all 346 rows of subset A,
each forecast made with its campaign held out. (a)~The duration curve predicts nearly one value per scan day. (b)~\Jev{}
follows the between-muscle spread, which is where its gain comes from; both regress towards the mean at the extremes.
Dashed line: identity.}
\label{fig:jevscatter}
\end{figure}

\subsection{Tier 3: is the gain signal or recognition?}
\label{sec:results-ablation}

Because more than half of the gain came from three much-published campaigns, the ablations and the recognition probe
were declared and run (1,280 new answers). \Cref{tab:ablation} and panel~a of \cref{fig:jevchecks} report them.

\begin{table}[htbp]
\centering
\caption[Tier 3: ablations]{Tier 3 ablations on the arm without history (346 rows, 32 campaigns). The paired gain is
over the duration curve; the change is relative to the full run (negative = worse). Intervals are campaign bootstraps.}
\label{tab:ablation}
\footnotesize
\begin{tabular*}{\linewidth}{@{\extracolsep{\fill}}l c c c c c c@{}}
\toprule
Variant & MAE & CRPS & \makecell{80\%\\ cover} & \makecell{Paired gain,\\ pp (95\% CI)} & \makecell{Change vs.\ full,\\ pp (95\% CI)} & Closer \\
\midrule
\inputtable{generated/tab_ablation.tex}
\bottomrule
\end{tabular*}
\end{table}

\paragraph{Against the reading rules}
\begin{itemize}
  \item \textbf{The generic variant keeps the gain}: 0.30\pp{} (0.01 to 0.62), at least half of the full run's 0.42\pp{}
        with an interval above zero; its difference from the full run, $-0.11$\pp, has an interval that includes zero.
        The gain does not need anything that identifies a study.
  \item \textbf{The model reads the reference data}: with the other campaigns' values shuffled, the gain becomes a loss
        of 1.25\pp; without any reference data the model is far worse than the curve (5.85 against 3.13\pp).
  \item \textbf{Recognition is not a live explanation} (\cref{fig:recognition}). The model places a mean probability
        of 0.14 on the true campaign, against 0.08 by chance, and names it first on 19\% of rows. Three campaigns pass
        the recognition bar---NASA SPRINT and AGBRESA (named first on 67\% of their rows) and VBR (50\%)---but of the
        three that carry most of the gain only SPRINT is recognised; for BBR2-2 and MEDES LTBR the model's commonest
        answer is ``none of these''. Recognition and gain correlate at $\rho = 0.37$ ($p = 0.26$, 11 campaigns).
\end{itemize}
By the declared table this is the first case: \emph{tier 3's point forecasts may be quoted as a result}, with the
caveats of \cref{sec:results-tier3}. The campaign-level errors point the same way. A model recalling published numbers
should do best without reference data on the campaigns it can name; it does not, being worse than the curve on
AGBRESA (6.13 against 2.95\pp) and VBR (5.73 against 0.95\pp). SPRINT is the exception (4.80 against 5.92\pp) and the
one place recognition may contribute, but the generic variant, which gives the model nothing to recognise SPRINT by,
still beats the curve there (3.79 against 5.92\pp). The one campaign no paper reports, UTMB Campaign~3, computed from
NASA's raw archive, also favours the model (0.67\pp{} gain, four rows).

\begin{figure}[htbp]
\centering
\includegraphics[width=0.7\linewidth]{fig_recognition}
\caption[Tier 3: recognition probe]{Recognition probe, 11 campaigns with proper names. Horizontal: mean probability
the model places on the true campaign name when shown the description the forecast uses; vertical: the campaign's
paired gain over the curve. Marker area is proportional to rows. Two of the three campaigns carrying most of the gain
(BBR2-2, MEDES LTBR) are not recognised.}
\label{fig:recognition}
\end{figure}

\subsection{Tier 3: does it hold up?}
\label{sec:results-validation}

The validation battery added 1,195 answers and 20 repeats past the cache (\cref{tab:validation}, \cref{fig:jevchecks}).

\begin{table}[htbp]
\centering
\caption[Tier 3: validation battery]{Tier 3 validation battery, leave-one-campaign-out throughout. MAE in pp for
\Jev{} and for the arm's reference; paired gain over the reference and change from the full run with 95\%
campaign-bootstrap intervals. On the history arm each interval rests on five campaigns.}
\label{tab:validation}
\footnotesize
\begin{tabular*}{\linewidth}{@{\extracolsep{\fill}}l c c c c@{}}
\toprule
Run & \makecell{\Jev{}\\ MAE} & \makecell{Reference\\ MAE} & \makecell{Paired gain,\\ pp (95\% CI)} & \makecell{Change vs.\ full,\\ pp (95\% CI)} \\
\midrule
\inputtable{generated/tab_validation.tex}
\bottomrule
\end{tabular*}
\end{table}

\begin{figure}[htbp]
\centering
\includegraphics[width=\linewidth]{fig_jev_checks}
\caption[Tier 3: the gain under every check]{Paired MAE gain of \Jev{} over each arm's reference under every
ablation and presentation check, with 95\% campaign-bootstrap intervals. Blue: interval above zero; orange: below
zero; grey: spans zero. (a)~Without history the gain survives every change of presentation and the removal of
identifying detail, and collapses when the reference data are shuffled or removed. (b)~With history the gain is
unchanged when the campaign's own earlier scans are shuffled.}
\label{fig:jevchecks}
\end{figure}

\paragraph{Against the reading rules}
\begin{itemize}
  \item \textbf{Robust to presentation: yes, on both arms.} Without history the MAE moves by 3\% under reordering and
        1\% under shifted ranges, and the gain keeps an interval above zero under both (0.34 and 0.45\pp). On the
        history arm both moves are 2\% and the gain keeps its sign.
  \item \textbf{Consistent: no, by the letter of the rule.} None of the 20 repeated requests returned identical
        probabilities. The probabilities moved by 0.034 on average (0.09 at most), the forecasts by 0.17\pp{} on
        average (0.48 at most), the most likely range held in 18 of 20, and the error on the repeated rows went from
        2.456 to 2.479\pp. Run-to-run variation is of the order of the differences between presentation runs and far
        below the 0.42\pp{} gain.
  \item \textbf{Reads the earlier scans: no.} Shuffling the held-out campaign's earlier values leaves the error where
        it was ($+0.06$\pp, interval across zero).
  \item \textbf{Reads the other campaigns on the history arm: not shown.} Both reference ablations are worse in their
        point estimates, by 0.13 and 0.52\pp, but on five campaigns neither interval clears zero.
  \item \textbf{The earlier scans help---strongly, but not through the model.} On the same 84 rows the model's error
        falls from 5.83 to 2.84\pp{} when it forecasts the change since the last scan rather than the level (paired
        2.98\pp, 1.54 to 5.51), and the baselines' from 7.27 to 3.21\pp{} (4.06\pp, 2.52 to 5.99). What the earlier
        scans supply is the anchor, and the code supplies it by adding the forecast change to the last scan.
\end{itemize}
Of the three conditions for calling tier 3 working rather than lucky, the one that carries the quoted number holds:
the result without history survives both presentation checks and its run-to-run variation is far smaller than the
gain. The other two fail and are reported as failing. The tier-3 number is therefore quoted from the arm without
history, with the note that repeated runs differ slightly; nothing is claimed about the model using a campaign's
earlier scans.

\subsection{The effect of correcting the time axis}
\label{sec:results-correction}

\Cref{tab:correction} records what changed when tiers 1 and 2 moved from the campaign's planned length to the day of
the scan (\cref{sec:target}). The fit improved by about 65 AIC points on identical rows, the plainest evidence that
the scan day is the right exposure; the time constant moved from the edge of the declared grid to its interior; and
the share of variance between muscles rose. What matters for the conclusions did not change: the duration forms remain
indistinguishable, the ranking keeps its order, and no tier-2 family beats the curve. Tier 3 was fitted on the scan
day from the start.

\begin{table}[htbp]
\centering
\caption[The time-axis correction]{Results fitted on the campaign's planned length (first runs, superseded) and on the
day of the scan (reported). Values for the planned length are as recorded in \file{framework/DESIGN.md}, section 9.4.}
\label{tab:correction}
\small
\begin{tabularx}{\linewidth}{@{}Y c c@{}}
\toprule
Result & Planned length & Day of the scan \\
\midrule
Tier 1, saturating form, AIC on the same 346 rows & 2152.3 & \textbf{2085.9} \\
Residual variance & 17.7 & \textbf{13.8} \\
Loss per doubling of days, logarithmic form (pp) & $-2.41$ ($-3.26$ to $-1.56$) & \textbf{$-$2.99} ($-3.44$ to $-2.53$) \\
Saturating time constant & 90 days, at the edge of the grid & \textbf{60 days}, inside it \\
Eventual loss, saturating form (\%) & $-19.9$ ($-24.8$ to $-15.0$) & \textbf{$-$17.3} ($-20.9$ to $-13.7$) \\
Share of variance between muscles within a campaign & 23\% & \textbf{30\%} \\
Baseline (logarithmic), out-of-campaign MAE & 3.21\pp & \textbf{3.16\pp} \\
Baseline pooled $R^2$ & 0.03 & \textbf{0.14} \\
Best tier-2 family against the baseline & SVR, $-0.4\%$ & \textbf{Random forest, $-$0.7\%} \\
Tier-2 pooled $R^2$ & 0.10--0.22 & \textbf{0.27--0.36} \\
Features stable in the top three ($\geq$80\% of folds) & duration only, at 69\% & \textbf{duration (91\%), plantar flexors (88\%)} \\
\bottomrule
\end{tabularx}
\end{table}

\subsection{Assessment against the declared criteria}
\label{sec:results-criteria}

\Cref{tab:assessment} sets every criterion declared before modelling, and every bar declared for tier 3, against
its outcome. Of the six criteria for tiers 1 and 2, the three that were allowed to fail split one met and two missed;
of the three blocking criteria, two are met and the physiological sign-off is outstanding.

\begin{table}[H]
\centering
\caption[Assessment against the declared criteria]{Assessment against the criteria of \cref{tab:criteria} and the
tier-3 bars.}
\label{tab:assessment}
\small
\begin{tabularx}{\linewidth}{@{}L{4.8cm} L{2.2cm} Y@{}}
\toprule
Criterion & Outcome & Evidence \\
\midrule
Out-of-campaign MAE $\leq$ 4\pp & Met & Every baseline form and every family, 3.16--3.40\pp \\
Best model $\geq$ 15\% better than baseline & Not met & Best family $-0.7\%$; null result reported \\
\quad TabPFN, post hoc, against the same bar & Not met & $-2.1\%$; pooled $R^2$ 0.39 \\
SHAP top-3 stability $\geq$ 80\% & Not met & Top two stable (91\%, 88\%); third place 59\% \\
Leakage audit & Met & Disjointness asserted in every fold; the assertion is tested against a bad split \\
Reproducibility & Met & One command rebuilds every result; 234 tests pass; tier 3 rebuilt from its cache \\
Physiological sign-off & \textbf{Outstanding} & Muscle classification and fitted relationships await the scientific lead \\
\addlinespace
Tier 3, MAE $\geq$ 15\% better & Not met & 13.3\% without history, 11.5\% with \\
Tier 3, CRPS $\geq$ 15\% better & Met without history & 17.7\%; with history 2.2\% worse \\
Tier 3, 80\% coverage in 70--90\% & Not met & 62\% and 38\%: ranges overconfident \\
Tier 3, leakage and reproducibility & Met & Assertions on every request; cache rebuild without API calls \\
\bottomrule
\end{tabularx}
\end{table}
